\documentclass[12pt]{article}
\usepackage{amsmath, amssymb, amsfonts}
\usepackage{geometry}
\geometry{margin=1in}

\title{Lecture 20: Linear Transformations and Expectations of Multiple RVs}
\date{}

\begin{document}

\maketitle

\section{Outline of the Lecture}

The lecture covers:
\begin{itemize}
\item Statistical Representation of Random Vectors
\item Mean Vector, Correlation, and Covariance Matrices
\item Linear Transformations
\item Mapping properties under $Y = AX + b$
\item Multivariate Gaussian Distributions
\item From 1D/2D cases to N-dimensional Joint PDFs
\item Special Cases and Properties
\item Implications of zero correlation and matrix factorization
\end{itemize}

\section{Random Vector Representation}

\subsection{Vector Form of Multiple Random Variables}

\textbf{Definition: Random Vector}

For random variables $X_1, X_2, \ldots, X_N$, we represent them as a vector:

\[
X =
\begin{bmatrix}
X_1 \\
X_2 \\
\vdots \\
X_N
\end{bmatrix}
\]

\textbf{Statements}
\begin{itemize}
\item $X$ is called a random vector.
\item It provides a compact representation of multiple random variables.
\end{itemize}

\section{Mean Vector}

\subsection{Expectation of a Random Vector}

\textbf{Mean Vector Definition}

\[
\mu_X = E[X]
\]

\textbf{Statements}
\begin{itemize}
\item Represents the average value of the random vector.
\item Each element corresponds to the mean of one random variable.
\end{itemize}

\subsection{Example Structure}

\[
\mu_X =
\begin{bmatrix}
E[X_1] \\
E[X_2] \\
\vdots \\
E[X_N]
\end{bmatrix}
\]

\section{Correlation Matrix}

\subsection{Correlation Between Multiple RVs}

\textbf{Correlation Matrix Definition}

\[
R_{XX} = E[XX^T]
\]

with:

\[
R_{XX}(i,j) = E[X_i X_j], \quad i, j \leq N
\]

\subsection{Matrix Representation}

\[
R_{XX} =
\begin{bmatrix}
E[X_1X_1] & E[X_1X_2] & \cdots & E[X_1X_N] \\
E[X_2X_1] & E[X_2X_2] & \cdots & E[X_2X_N] \\
\vdots & \vdots & \ddots & \vdots \\
E[X_NX_1] & E[X_NX_2] & \cdots & E[X_NX_N]
\end{bmatrix}
\]

\textbf{Property}

$R_{XX}$ is a symmetric matrix.

\section{Covariance Matrix}

\subsection{Mathematical Definition}

\[
C_{XX} = E[(X - \mu_X)(X - \mu_X)^T]
\]

\subsection{Practical Structure}

\[
C_{XX} =
\begin{bmatrix}
\text{Var}(X_1) & \text{Cov}(X_1, X_2) & \cdots \\
\text{Cov}(X_2, X_1) & \text{Var}(X_2) & \cdots \\
\vdots & \vdots & \ddots
\end{bmatrix}
\]

\subsection{Symmetry Property}

\[
\text{Cov}(X_i, X_j) = \text{Cov}(X_j, X_i)
\]

The matrix is a “mirror” of itself.

\section{Linear Transformation of Random Vectors}

\subsection{The Linear Model}

\[
Y = AX + b
\]

\textbf{Definitions}
\begin{itemize}
\item $A$: Transformation matrix (scales and rotates the data)
\item $b$: Constant vector (shifts/translates the center)
\item $X$: Input random vector
\end{itemize}

\subsection{Expanded Form}

\begin{align*}
Y_1 &= a_{11}X_1 + a_{12}X_2 + \cdots + a_{1N}X_N + b_1 \\
Y_2 &= a_{21}X_1 + a_{22}X_2 + \cdots + a_{2N}X_N + b_2 \\
\vdots \\
Y_N &= a_{N1}X_1 + a_{N2}X_2 + \cdots + a_{NN}X_N + b_N
\end{align*}

\section{Mean of Transformed Vector}

\subsection{Definition}

\[
\mu_Y = E[Y]
\]

\subsection{Derivation (exact steps)}

\begin{align*}
\mu_Y &= E[AX + b] \\
\mu_Y &= E[AX] + E[b]
\end{align*}

\subsection{Result}

\[
\mu_Y = A\mu_X + b
\]

\section{Correlation Matrix Transformation}

\subsection{Definition}

\[
R_{YY} = E[YY^T] = E[(AX + b)(AX + b)^T]
\]

\subsection{Expansion of Terms}

\[
R_{YY} = E[AXX^T A^T + AXb^T + bX^T A^T + bb^T]
\]

\subsection{Applying Expectation}

\[
R_{YY} = A E[XX^T] A^T + A E[X] b^T + b E[X]^T A^T + bb^T
\]

\subsection{Final Result}

\[
R_{YY} = A R_{XX} A^T + A\mu_X b^T + b\mu_X^T A^T + bb^T
\]

\section{Covariance Matrix Transformation}

\subsection{Key Observation}

\[
Y - \mu_Y = (AX + b) - (A\mu_X + b) = A(X - \mu_X)
\]

The deviation from the mean is independent of the shift $b$.

\subsection{Derivation}

\begin{align*}
C_{YY} &= E[(Y - \mu_Y)(Y - \mu_Y)^T] \\
C_{YY} &= E[A(X - \mu_X)(X - \mu_X)^T A^T]
\end{align*}

\subsection{Simplified Result}

\[
C_{YY} = A C_{XX} A^T
\]

The shift $b$ does not affect the covariance (spread).

\section{Summary: Effects of Linear Transformation}

If $Y = AX + b$:

\[
\mu_Y = A\mu_X + b
\]

\[
C_{YY} = A C_{XX} A^T
\]

\[
R_{YY} = A R_{XX} A^T + A\mu_X b^T + b\mu_X^T A^T + bb^T
\]

\section{Example: Numerical Linear Transformation}

\subsection{Given}

\[
\mu_X =
\begin{bmatrix}
2 \\
-1 \\
1
\end{bmatrix}, \quad
A =
\begin{bmatrix}
1 & 0 & -1 \\
0 & -1 & 1 \\
-1 & 1 & 0
\end{bmatrix}
\]

\[
R_{XX} =
\begin{bmatrix}
13 & 2 & 3 \\
2 & 10 & 3 \\
3 & 3 & 10
\end{bmatrix}
\]

\subsection{(a) Mean Vector}

\[
\mu_Y = A\mu_X
\]

\[
\mu_Y =
\begin{bmatrix}
(1)(2) + (0)(-1) + (-1)(1) \\
(0)(2) + (-1)(-1) + (1)(1) \\
(-1)(2) + (1)(-1) + (0)(1)
\end{bmatrix}
\]

\[
\mu_Y =
\begin{bmatrix}
1 \\
2 \\
-3
\end{bmatrix}
\]

\subsection{(b) Correlation Matrix}

\[
R_{YY} = A R_{XX} A^T
\]

\[
AR_{XX} =
\begin{bmatrix}
10 & -1 & -7 \\
1 & -7 & 7 \\
-11 & 8 & 0
\end{bmatrix}
\]

\[
R_{YY} =
\begin{bmatrix}
10 & -1 & -7 \\
1 & -7 & 7 \\
-11 & 8 & 0
\end{bmatrix}
\begin{bmatrix}
1 & 0 & -1 \\
0 & -1 & 1 \\
-1 & 1 & 0
\end{bmatrix}
\]

\[
R_{YY} =
\begin{bmatrix}
17 & -6 & -11 \\
-6 & 14 & -8 \\
-11 & -8 & 19
\end{bmatrix}
\]

\subsection{(c) Covariance Matrix}

\[
C_{YY} = R_{YY} - \mu_Y \mu_Y^T
\]

\[
\mu_Y \mu_Y^T =
\begin{bmatrix}
1 & 2 & -3 \\
2 & 4 & -6 \\
-3 & -6 & 9
\end{bmatrix}
\]

\[
C_{YY} =
\begin{bmatrix}
17 & -6 & -9 \\
-6 & 14 & -8 \\
-9 & -8 & 19
\end{bmatrix}
-
\begin{bmatrix}
1 & 2 & -3 \\
2 & 4 & -6 \\
-3 & -6 & 9
\end{bmatrix}
\]

\[
C_{YY} =
\begin{bmatrix}
16 & -8 & -8 \\
-8 & 10 & -2 \\
-8 & -2 & 10
\end{bmatrix}
\]

\section{Gaussian Distributions}

\subsection{Univariate Gaussian (1D)}

\[
f_X(x) =
\frac{1}{\sqrt{2\pi\sigma^2}}
\exp\left(-\frac{(x-\mu)^2}{2\sigma^2}\right)
\]

\subsection{Joint Gaussian (2D)}

For two RVs $X$ and $Y$, the joint PDF is given.

\section{Multivariate Gaussian Distribution}

\subsection{Definition}

\[
X = [X_1, X_2, \ldots, X_N]^T
\]

\subsection{General Vector PDF}

\[
f_X(x) =
\frac{1}{\sqrt{(2\pi)^N \det(C_{XX})}}
\exp\left(-\frac{1}{2}(x-\mu_X)^T C_{XX}^{-1}(x-\mu_X)\right)
\]

\subsection{Structure Components}

\begin{itemize}
\item $\mu_X$: Mean vector $(N \times 1)$
\item $C_{XX}$: Covariance matrix $(N \times N)$
\end{itemize}

\section{Example: 2D Joint Gaussian PDF}

\subsection{Given}

\[
\mu_X =
\begin{bmatrix}
\mu_1 \\
\mu_2
\end{bmatrix}
\]

\[
C_{XX} =
\begin{bmatrix}
\sigma_1^2 & \rho\sigma_1\sigma_2 \\
\rho\sigma_1\sigma_2 & \sigma_2^2
\end{bmatrix}
\]

\subsection{Determinant}

\[
\det(C_{XX}) = \sigma_1^2 \sigma_2^2 (1 - \rho^2)
\]

\subsection{Inverse}

\[
C_{XX}^{-1} = \frac{1}{\det(C_{XX})} \text{adj}(C_{XX})
\]

\subsection{Quadratic Form}

\[
Q(x) = (x - \mu_X)^T C_{XX}^{-1}(x - \mu_X)
\]

\subsection{Final Expression}

\[
Q(x) =
\frac{1}{1-\rho^2}
\left[
\left(\frac{x_1-\mu_1}{\sigma_1}\right)^2
-2\rho\left(\frac{x_1-\mu_1}{\sigma_1}\frac{x_2-\mu_2}{\sigma_2}\right)
+\left(\frac{x_2-\mu_2}{\sigma_2}\right)^2
\right]
\]

\section{Special Case: Uncorrelated Gaussian Variables}

\subsection{Key Insight}

\[
\text{Cov}(X_i, X_j) = 0, \quad \forall i \neq j
\]

\subsection{Covariance Matrix Simplification}

Diagonal matrix form.

\subsection{Result}

\[
f_X(x) = \prod_{n=1}^{N}
\frac{1}{\sqrt{2\pi\sigma_n^2}}
\exp\left(-\frac{(x_n-\mu_n)^2}{2\sigma_n^2}\right)
\]

\subsection{Interpretation}

\begin{itemize}
\item Covariance matrix is diagonal
\item Joint Gaussian PDF factorizes into product of individual PDFs
\end{itemize}

\section{Homework}

\[
\text{Var}[X+Y] = \text{Var}[X] + \text{Var}[Y] + 2\text{Cov}(X,Y)
\]

\section{Case Study: Smart City Traffic + Weather System}

\subsection{Step 1: Define Random Vector}

\[
X =
\begin{bmatrix}
X_1 \\
X_2 \\
X_3 \\
X_4 \\
X_5
\end{bmatrix}
\]

\subsection{Step 2: Mean Vector}

\[
\mu = E[X]
\]

Example values:
\begin{itemize}
\item Traffic = 120
\item Speed = 35
\item Rainfall = 2
\item Temperature = 32
\item AQI = 180
\end{itemize}

\subsection{Step 3: Covariance Matrix}

Examples:
\begin{itemize}
\item $\text{Cov}(X_1, X_3) > 0$
\item $\text{Cov}(X_2, X_3) < 0$
\item $\text{Cov}(X_1, X_5) > 0$
\item $\text{Cov}(X_4, X_5) < 0$
\end{itemize}

Covariance matrix = dependency map of the city.

\subsection{Step 4: Joint Gaussian Distribution}

\[
X \sim \mathcal{N}(\mu, \Sigma)
\]

\subsection{Step 6: Transformation}

\[
Y = 0.6X_1 - 0.4X_2 + 0.8X_3
\]

\[
Y = a^T X
\]

\subsection{Inference}

If rainfall increases suddenly:
\begin{itemize}
\item Speed $\downarrow$
\item Traffic $\uparrow$
\item Pollution $\uparrow$
\end{itemize}

\end{document}